% Options for packages loaded elsewhere
\PassOptionsToPackage{unicode}{hyperref}
\PassOptionsToPackage{hyphens}{url}
%
\documentclass[
]{article}
\usepackage{amsmath,amssymb}
\usepackage{iftex}
\ifPDFTeX
  \usepackage[T1]{fontenc}
  \usepackage[utf8]{inputenc}
  \usepackage{textcomp} % provide euro and other symbols
\else % if luatex or xetex
  \usepackage{unicode-math} % this also loads fontspec
  \defaultfontfeatures{Scale=MatchLowercase}
  \defaultfontfeatures[\rmfamily]{Ligatures=TeX,Scale=1}
\fi
\usepackage{lmodern}
\ifPDFTeX\else
  % xetex/luatex font selection
\fi
% Use upquote if available, for straight quotes in verbatim environments
\IfFileExists{upquote.sty}{\usepackage{upquote}}{}
\IfFileExists{microtype.sty}{% use microtype if available
  \usepackage[]{microtype}
  \UseMicrotypeSet[protrusion]{basicmath} % disable protrusion for tt fonts
}{}
\makeatletter
\@ifundefined{KOMAClassName}{% if non-KOMA class
  \IfFileExists{parskip.sty}{%
    \usepackage{parskip}
  }{% else
    \setlength{\parindent}{0pt}
    \setlength{\parskip}{6pt plus 2pt minus 1pt}}
}{% if KOMA class
  \KOMAoptions{parskip=half}}
\makeatother
\usepackage{xcolor}
\usepackage[margin=1in]{geometry}
\usepackage{color}
\usepackage{fancyvrb}
\newcommand{\VerbBar}{|}
\newcommand{\VERB}{\Verb[commandchars=\\\{\}]}
\DefineVerbatimEnvironment{Highlighting}{Verbatim}{commandchars=\\\{\}}
% Add ',fontsize=\small' for more characters per line
\usepackage{framed}
\definecolor{shadecolor}{RGB}{248,248,248}
\newenvironment{Shaded}{\begin{snugshade}}{\end{snugshade}}
\newcommand{\AlertTok}[1]{\textcolor[rgb]{0.94,0.16,0.16}{#1}}
\newcommand{\AnnotationTok}[1]{\textcolor[rgb]{0.56,0.35,0.01}{\textbf{\textit{#1}}}}
\newcommand{\AttributeTok}[1]{\textcolor[rgb]{0.13,0.29,0.53}{#1}}
\newcommand{\BaseNTok}[1]{\textcolor[rgb]{0.00,0.00,0.81}{#1}}
\newcommand{\BuiltInTok}[1]{#1}
\newcommand{\CharTok}[1]{\textcolor[rgb]{0.31,0.60,0.02}{#1}}
\newcommand{\CommentTok}[1]{\textcolor[rgb]{0.56,0.35,0.01}{\textit{#1}}}
\newcommand{\CommentVarTok}[1]{\textcolor[rgb]{0.56,0.35,0.01}{\textbf{\textit{#1}}}}
\newcommand{\ConstantTok}[1]{\textcolor[rgb]{0.56,0.35,0.01}{#1}}
\newcommand{\ControlFlowTok}[1]{\textcolor[rgb]{0.13,0.29,0.53}{\textbf{#1}}}
\newcommand{\DataTypeTok}[1]{\textcolor[rgb]{0.13,0.29,0.53}{#1}}
\newcommand{\DecValTok}[1]{\textcolor[rgb]{0.00,0.00,0.81}{#1}}
\newcommand{\DocumentationTok}[1]{\textcolor[rgb]{0.56,0.35,0.01}{\textbf{\textit{#1}}}}
\newcommand{\ErrorTok}[1]{\textcolor[rgb]{0.64,0.00,0.00}{\textbf{#1}}}
\newcommand{\ExtensionTok}[1]{#1}
\newcommand{\FloatTok}[1]{\textcolor[rgb]{0.00,0.00,0.81}{#1}}
\newcommand{\FunctionTok}[1]{\textcolor[rgb]{0.13,0.29,0.53}{\textbf{#1}}}
\newcommand{\ImportTok}[1]{#1}
\newcommand{\InformationTok}[1]{\textcolor[rgb]{0.56,0.35,0.01}{\textbf{\textit{#1}}}}
\newcommand{\KeywordTok}[1]{\textcolor[rgb]{0.13,0.29,0.53}{\textbf{#1}}}
\newcommand{\NormalTok}[1]{#1}
\newcommand{\OperatorTok}[1]{\textcolor[rgb]{0.81,0.36,0.00}{\textbf{#1}}}
\newcommand{\OtherTok}[1]{\textcolor[rgb]{0.56,0.35,0.01}{#1}}
\newcommand{\PreprocessorTok}[1]{\textcolor[rgb]{0.56,0.35,0.01}{\textit{#1}}}
\newcommand{\RegionMarkerTok}[1]{#1}
\newcommand{\SpecialCharTok}[1]{\textcolor[rgb]{0.81,0.36,0.00}{\textbf{#1}}}
\newcommand{\SpecialStringTok}[1]{\textcolor[rgb]{0.31,0.60,0.02}{#1}}
\newcommand{\StringTok}[1]{\textcolor[rgb]{0.31,0.60,0.02}{#1}}
\newcommand{\VariableTok}[1]{\textcolor[rgb]{0.00,0.00,0.00}{#1}}
\newcommand{\VerbatimStringTok}[1]{\textcolor[rgb]{0.31,0.60,0.02}{#1}}
\newcommand{\WarningTok}[1]{\textcolor[rgb]{0.56,0.35,0.01}{\textbf{\textit{#1}}}}
\usepackage{graphicx}
\makeatletter
\newsavebox\pandoc@box
\newcommand*\pandocbounded[1]{% scales image to fit in text height/width
  \sbox\pandoc@box{#1}%
  \Gscale@div\@tempa{\textheight}{\dimexpr\ht\pandoc@box+\dp\pandoc@box\relax}%
  \Gscale@div\@tempb{\linewidth}{\wd\pandoc@box}%
  \ifdim\@tempb\p@<\@tempa\p@\let\@tempa\@tempb\fi% select the smaller of both
  \ifdim\@tempa\p@<\p@\scalebox{\@tempa}{\usebox\pandoc@box}%
  \else\usebox{\pandoc@box}%
  \fi%
}
% Set default figure placement to htbp
\def\fps@figure{htbp}
\makeatother
\setlength{\emergencystretch}{3em} % prevent overfull lines
\providecommand{\tightlist}{%
  \setlength{\itemsep}{0pt}\setlength{\parskip}{0pt}}
\setcounter{secnumdepth}{-\maxdimen} % remove section numbering
\usepackage{bookmark}
\IfFileExists{xurl.sty}{\usepackage{xurl}}{} % add URL line breaks if available
\urlstyle{same}
\hypersetup{
  pdftitle={ProblemSet\_1},
  pdfauthor={Zachary Johnigan},
  hidelinks,
  pdfcreator={LaTeX via pandoc}}

\title{ProblemSet\_1}
\author{Zachary Johnigan}
\date{2026-07-22}

\begin{document}
\maketitle

Zachary Johnigan Aida Pacheco-Applegate Data For Policy Analys/Mgmt SSAD
48500/1 {[}60558{]} Due: Sun Aug 2, 2026 11:59pm

\centering

ProblemSet\_1

\begin{enumerate}
\def\labelenumi{\arabic{enumi}.}
\tightlist
\item
  Prepare the data for analysis. {[}x{]} Subset the NSECE dataset to
  include the following variables: o WF9\_WoRK\_WAGE (hourly wage) o
  WF9\_WoRK\_RoLE (staff role) o WF9\_WoRK\_YRS (number of years worked
  in the program) o WF9\_CHAR\_HISP (Hispanic/Latino ethnicity) o
  WF9\_CHAR\_RACE (race) o WF9\_CHAR\_EDUC (highest level of education
  completed) o WF9\_CHAR\_GENDER (gender)
\end{enumerate}

{[}x{]} Rename the variables using clear, descriptive names. {[}-{]}
Perform any data cleaning you consider necessary. {[}x{]} For each
variable, identify its level of measurement (nominal, ordinal, or
interval/ratio) and whether it is discrete or continuous. {[}-{]}
Include the appropriate descriptive statistics, tables, or graphs, as
well as measures of central tendency, dispersion, and shape when
applicable.

Describe and interpret your findings. (4 points)

\section{1. Set up working space and import
packages}\label{set-up-working-space-and-import-packages}

\begin{Shaded}
\begin{Highlighting}[]
\CommentTok{\# This cleans up everything in the "Environment" pane.}
\FunctionTok{rm}\NormalTok{(}\AttributeTok{list=}\FunctionTok{ls}\NormalTok{())}
\FunctionTok{options}\NormalTok{(}\StringTok{"error"}\NormalTok{)}
\end{Highlighting}
\end{Shaded}

\begin{verbatim}
## $error
## NULL
\end{verbatim}

\begin{Shaded}
\begin{Highlighting}[]
\FunctionTok{options}\NormalTok{(}\AttributeTok{lifecycle\_disable\_verbose\_retirement =} \ConstantTok{TRUE}\NormalTok{)}
\end{Highlighting}
\end{Shaded}

\begin{Shaded}
\begin{Highlighting}[]
\CommentTok{\# Import packages}
\NormalTok{package.list }\OtherTok{\textless{}{-}} \FunctionTok{c}\NormalTok{(}\StringTok{"tidyverse"}\NormalTok{)}

\CommentTok{\# read the \textasciigrave{}package.list\textasciigrave{} for the list of packages we will be needing}
\CommentTok{\# Check what packages are installed already}
\CommentTok{\# If any of the packages in the c() list are not currently installed, install them.}
\CommentTok{\# Note: Check the "Packages" tab to verify that they were installed. If they were not, then check the box and click "Install".}
\ControlFlowTok{for}\NormalTok{ (p }\ControlFlowTok{in}\NormalTok{ package.list)\{}
  \ControlFlowTok{if}\NormalTok{ (}\SpecialCharTok{!}\NormalTok{p }\SpecialCharTok{\%in\%} \FunctionTok{installed.packages}\NormalTok{()[, }\StringTok{"Package"}\NormalTok{]) }\FunctionTok{install.packages}\NormalTok{(p)}
  \FunctionTok{library}\NormalTok{(p, }\AttributeTok{character.only=}\ConstantTok{TRUE}\NormalTok{)}
\NormalTok{\}}
\end{Highlighting}
\end{Shaded}

\begin{verbatim}
## -- Attaching core tidyverse packages ------------------------ tidyverse 2.0.0 --
## v dplyr     1.2.1     v readr     2.2.0
## v forcats   1.0.1     v stringr   1.6.0
## v ggplot2   4.0.3     v tibble    3.3.1
## v lubridate 1.9.5     v tidyr     1.3.2
## v purrr     1.2.2     
## -- Conflicts ------------------------------------------ tidyverse_conflicts() --
## x dplyr::filter() masks stats::filter()
## x dplyr::lag()    masks stats::lag()
## i Use the conflicted package (<http://conflicted.r-lib.org/>) to force all conflicts to become errors
\end{verbatim}

\begin{Shaded}
\begin{Highlighting}[]
\CommentTok{\# Set up working directory}

\DocumentationTok{\#\# check your directory}
\FunctionTok{getwd}\NormalTok{()}
\end{Highlighting}
\end{Shaded}

\begin{verbatim}
## [1] "/home/zacharyjohnigan/UChicago_Student/2_Data_For_Policy_Analys_-_Mgmt/Problem_Sets/Problem_Set_1"
\end{verbatim}

\begin{Shaded}
\begin{Highlighting}[]
\DocumentationTok{\#\# set up working directory}
\NormalTok{my\_directory   }\OtherTok{\textless{}{-}} \StringTok{"/home/zacharyjohnigan/UChicago\_Student/2\_Data\_For\_Policy\_Analys\_{-}\_Mgmt/"}
\NormalTok{my\_directory}
\end{Highlighting}
\end{Shaded}

\begin{verbatim}
## [1] "/home/zacharyjohnigan/UChicago_Student/2_Data_For_Policy_Analys_-_Mgmt/"
\end{verbatim}

\begin{Shaded}
\begin{Highlighting}[]
\NormalTok{my\_data        }\OtherTok{\textless{}{-}} \FunctionTok{paste0}\NormalTok{(my\_directory, }\StringTok{"Data/"}\NormalTok{)}
\NormalTok{my\_data}
\end{Highlighting}
\end{Shaded}

\begin{verbatim}
## [1] "/home/zacharyjohnigan/UChicago_Student/2_Data_For_Policy_Analys_-_Mgmt/Data/"
\end{verbatim}

\begin{Shaded}
\begin{Highlighting}[]
\NormalTok{output\_dir     }\OtherTok{\textless{}{-}} \FunctionTok{paste0}\NormalTok{(my\_directory, }\StringTok{"Problem\_Sets/Problem\_Set\_1/"}\NormalTok{)}
\NormalTok{output\_dir}
\end{Highlighting}
\end{Shaded}

\begin{verbatim}
## [1] "/home/zacharyjohnigan/UChicago_Student/2_Data_For_Policy_Analys_-_Mgmt/Problem_Sets/Problem_Set_1/"
\end{verbatim}

\section{2. Load dataset}\label{load-dataset}

\begin{Shaded}
\begin{Highlighting}[]
\CommentTok{\# Import NSECE 2019 workforce questionnnaire (since the data is already in R}
\CommentTok{\# format we just need to use the line 52 of this code to load it into our work}
\CommentTok{\# environment)}

\FunctionTok{load}\NormalTok{(}\FunctionTok{paste0}\NormalTok{(my\_data, }\StringTok{"37941{-}0005{-}Data.rda"}\NormalTok{))}
\end{Highlighting}
\end{Shaded}

\section{4. Basic cleaning of the
data}\label{basic-cleaning-of-the-data}

\begin{Shaded}
\begin{Highlighting}[]
\CommentTok{\# Rename dataset}
\NormalTok{nsece\_2019\_wf }\OtherTok{\textless{}{-}}\NormalTok{ da37941}\FloatTok{.0005}
\end{Highlighting}
\end{Shaded}

\section{Manipulation of the data}\label{manipulation-of-the-data}

\section{Subset the NSECE dataset to include the following
variables:}\label{subset-the-nsece-dataset-to-include-the-following-variables}

\section{Subset the NSECE dataset}\label{subset-the-nsece-dataset}

\begin{Shaded}
\begin{Highlighting}[]
\NormalTok{nsece\_2019\_wf\_subset }\OtherTok{\textless{}{-}}
\NormalTok{  nsece\_2019\_wf }\SpecialCharTok{\%\textgreater{}\%}
  \FunctionTok{select}\NormalTok{(}
\NormalTok{    WF9\_WORK\_WAGE,}
\NormalTok{    WF9\_WORK\_ROLE,}
\NormalTok{    WF9\_WORK\_YRS,}
\NormalTok{    WF9\_CHAR\_HISP,}
\NormalTok{    WF9\_CHAR\_RACE,}
\NormalTok{    WF9\_CHAR\_EDUC,}
\NormalTok{    WF9\_CHAR\_GENDER) }\SpecialCharTok{\%\textgreater{}\%}
  \FunctionTok{rename}\NormalTok{(}
    \AttributeTok{hourly\_wage =}\NormalTok{ WF9\_WORK\_WAGE,}
    \AttributeTok{staff\_role =}\NormalTok{ WF9\_WORK\_ROLE,}
    \AttributeTok{number\_of\_years\_worked\_in\_the\_program =}\NormalTok{ WF9\_WORK\_YRS,}
    \AttributeTok{Hispanic\_Latino\_ethnicity =}\NormalTok{ WF9\_CHAR\_HISP,}
    \AttributeTok{race =}\NormalTok{ WF9\_CHAR\_RACE,}
    \AttributeTok{highest\_level\_of\_education\_completed =}\NormalTok{ WF9\_CHAR\_EDUC,}
    \AttributeTok{gender =}\NormalTok{ WF9\_CHAR\_GENDER}
\NormalTok{  )}
\end{Highlighting}
\end{Shaded}

\begin{Shaded}
\begin{Highlighting}[]
\NormalTok{nsece\_2019\_wf\_recode }\OtherTok{\textless{}{-}}
\NormalTok{  nsece\_2019\_wf\_subset }\SpecialCharTok{\%\textgreater{}\%}
  \FunctionTok{mutate}\NormalTok{(}

    \AttributeTok{gender\_char =} \FunctionTok{as.character}\NormalTok{(gender),}
    \AttributeTok{gender\_char =} \FunctionTok{case\_when}\NormalTok{(}
\NormalTok{      gender }\SpecialCharTok{==} \StringTok{"(2) Female"} \SpecialCharTok{\textasciitilde{}} \StringTok{"Female"}\NormalTok{,}
\NormalTok{      gender }\SpecialCharTok{==} \StringTok{"(1) Male"} \SpecialCharTok{\textasciitilde{}} \StringTok{"Male"}
\NormalTok{    ),}

    \AttributeTok{staff\_role\_char =} \FunctionTok{as.character}\NormalTok{(staff\_role),}
    \AttributeTok{staff\_role\_char =} \FunctionTok{case\_when}\NormalTok{(}
\NormalTok{      staff\_role }\SpecialCharTok{==} \StringTok{"(1) Aide or assistant teacher"} \SpecialCharTok{\textasciitilde{}} \StringTok{"Aide/Assistant"}\NormalTok{,}
\NormalTok{      staff\_role }\SpecialCharTok{==} \StringTok{"(2) Teacher, instructor, or lead teacher"} \SpecialCharTok{\textasciitilde{}} \StringTok{"Teacher/Lead"}\NormalTok{,}
\NormalTok{      staff\_role }\SpecialCharTok{==} \StringTok{"(3) Other/undetermined"} \SpecialCharTok{\textasciitilde{}} \StringTok{"Other"}
\NormalTok{    )}
\NormalTok{  )}
\end{Highlighting}
\end{Shaded}

\section{Each variable's level of measurement (nominal, ordinal, or
interval/ratio) and whether it is discrete or
continuous.}\label{each-variables-level-of-measurement-nominal-ordinal-or-intervalratio-and-whether-it-is-discrete-or-continuous.}

According to the 37941-0005-User\_guide.pdf, ``manual provides
information on the 2019 National Survey of Early Care and Education
(NSECE) workforce public-use data files''(), and the lecture we can
define each of the 7 variables as follows.

WF9\_WORK\_WAGE - hourly\_wage - Hourly wage for Workforce respondent
(top-coded) Level of Measurement: Ratio, Continuous

WF9\_WORK\_ROLE - staff\_role - Role of Workforce respondent Level of
Measurement: Ordinal, Discreate

WF9\_WORK\_YRS - number\_of\_years\_worked\_in\_the\_program - Num of
years respondent worked in program Level of Measurement: Ratio,
Continuous

WF9\_CHAR\_HISP - Hispanic\_Latino\_ethnicity - Respondent is of
Hispanic/Latino descent Level of Measurement: Nominal, Discrete

WF9\_CHAR\_RACE - race - Respondent's race Level of Measurement:
Nominal, Discrete

WF9\_CHAR\_EDUC - highest\_level\_of\_education\_completed - Highest
education level obtained Level of Measurement: Ordinal, Discrete

WF9\_CHAR\_GENDER - gender - The respondent's self-reported gender Level
of Measurement: Nominal, Discrete

\section{Citations}\label{citations}

NSECE Project Team (National Opinion Research Center). (2019). National
Survey of Early Care and Education (NSECE), {[}United States{]}, 2019:
NSECE user's guide for workforce public-use data file (ICPSR 37941).
Inter-university Consortium for Political and Social Research.

Pacheco-Applegate, A. (2026). Data for Policy Analysis and Management:
Class 2---Data measurement and variables {[}PowerPoint slides{]}. Harris
School of Public Policy, University of Chicago.

\begin{enumerate}
\def\labelenumi{\arabic{enumi}.}
\setcounter{enumi}{1}
\tightlist
\item
  Examine the following relationships. For each analysis, include the
  appropriate statistical test to determine whether the observed
  differences or associations are statistically significant. Interpret
  your results in the context of the ECE workforce. (4 points)
\end{enumerate}

\section{My Answer}\label{my-answer}

From lecture 4, we learn that a bivariate analysis describes, compares,
and interprets associations between variables.
\pandocbounded{\includegraphics[keepaspectratio]{ProblemSet_1_files/1 Types of Bivariate Analysis.png}}

We also learn about statistical tests for determining whether the
observed differences or associations are statistically significant.
There are five tests we can use to perform these analysis, including:
Two-sample t-test, Two-sample z-test, ANOVA, Chi-square test and Paired
test.

I will specify which test I used for each comparison and what sort of
relationship they have---either positive, negative or null.

In a bivariate analysis, one variable defines the groups (explanatory
variable) ⟶ X and the other variable is the outcome we measure (response
variable) ⟶ Y

Take hourly\_wage and staff\_role for example. We can start by asking
first; among the respondents, does hourly wage differ by staff role?
This seems like an easy question to answer, but we can graph this to
show the relationship visually. Let's try a two-sample t-test.

\begin{enumerate}
\def\labelenumi{\alph{enumi}.}
\tightlist
\item
  Hourly\_wage and staff\_role
\end{enumerate}

\begin{Shaded}
\begin{Highlighting}[]
\CommentTok{\# a. Hourly wage and staff role}

\CommentTok{\# Research question:}
\CommentTok{\# Does hourly wage differ by staff role among ECE workers?}

\NormalTok{wage\_by\_role }\OtherTok{\textless{}{-}}
\NormalTok{  nsece\_2019\_wf\_recode }\SpecialCharTok{\%\textgreater{}\%}
  \FunctionTok{select}\NormalTok{(staff\_role, hourly\_wage) }\SpecialCharTok{\%\textgreater{}\%}          \CommentTok{\# keep only variables of interest}
  \FunctionTok{drop\_na}\NormalTok{(staff\_role, hourly\_wage) }\SpecialCharTok{\%\textgreater{}\%}         \CommentTok{\# remove missing values}
  \FunctionTok{group\_by}\NormalTok{(staff\_role) }\SpecialCharTok{\%\textgreater{}\%}                     \CommentTok{\# calculate statistics within each role}
  \FunctionTok{summarize}\NormalTok{(}
    \AttributeTok{n =} \FunctionTok{n}\NormalTok{(),}
    \AttributeTok{mean\_wage =} \FunctionTok{mean}\NormalTok{(hourly\_wage),}
    \AttributeTok{median\_wage =} \FunctionTok{median}\NormalTok{(hourly\_wage),}
    \AttributeTok{sd\_wage =} \FunctionTok{sd}\NormalTok{(hourly\_wage),}
    \AttributeTok{min\_wage =} \FunctionTok{min}\NormalTok{(hourly\_wage),}
    \AttributeTok{max\_wage =} \FunctionTok{max}\NormalTok{(hourly\_wage)}
\NormalTok{  )}

\NormalTok{wage\_by\_role}
\end{Highlighting}
\end{Shaded}

\begin{verbatim}
## # A tibble: 3 x 7
##   staff_role                   n mean_wage median_wage sd_wage min_wage max_wage
##   <fct>                    <int>     <dbl>       <dbl>   <dbl>    <dbl>    <dbl>
## 1 (1) Aide or assistant t~  1469      13.7        12      9.17     1.5     108. 
## 2 (2) Teacher, instructor~  2785      16.2        13.4   11.3      1.29    108. 
## 3 (3) Other/undetermined      16      18.4        15.5    8.75     9        42.1
\end{verbatim}

The summary statistics show that average hourly wage varies across staff
roles, with teachers, instructors, and lead teachers earning a higher
average hourly wage than aides or assistant teachers. A graph of this
comparison

\begin{Shaded}
\begin{Highlighting}[]
\NormalTok{wage\_by\_role\_plot }\OtherTok{\textless{}{-}}
\NormalTok{  nsece\_2019\_wf\_recode }\SpecialCharTok{\%\textgreater{}\%}
  \FunctionTok{drop\_na}\NormalTok{(staff\_role, hourly\_wage) }\SpecialCharTok{\%\textgreater{}\%}
  \FunctionTok{ggplot}\NormalTok{(}\FunctionTok{aes}\NormalTok{(}\AttributeTok{x =}\NormalTok{ staff\_role, }\AttributeTok{y =}\NormalTok{ hourly\_wage)) }\SpecialCharTok{+}
  \FunctionTok{geom\_boxplot}\NormalTok{()}

\NormalTok{wage\_by\_role\_plot}
\end{Highlighting}
\end{Shaded}

\pandocbounded{\includegraphics[keepaspectratio]{ProblemSet_1_files/ProblemSet_1_files/figure-latex/unnamed-chunk-9-1.pdf}}
The boxplot shows that teachers, instructors, and lead teachers
generally earn higher hourly wages than aides or assistant teachers,
although there is considerable variation and some outliers within each
staff role.

\subsection{b. Hourly wage and number of years worked in the
program}\label{b.-hourly-wage-and-number-of-years-worked-in-the-program}

Take \textbf{hourly\_wage} and
\textbf{number\_of\_years\_worked\_in\_the\_program} for example. We can
start by asking first; among the respondents, does hourly wage differ
depending on the number of years worked in the program? Because both
variables are continuous, a scatterplot is an appropriate way to
visualize the relationship.

\begin{Shaded}
\begin{Highlighting}[]
\CommentTok{\# b. Hourly wage and number of years worked in the program}

\CommentTok{\# Research question:}
\CommentTok{\# Does hourly wage differ by the number of years worked in the program among ECE workers?}

\NormalTok{wage\_by\_years }\OtherTok{\textless{}{-}}

\NormalTok{  nsece\_2019\_wf\_recode }\SpecialCharTok{\%\textgreater{}\%}

  \FunctionTok{select}\NormalTok{(number\_of\_years\_worked\_in\_the\_program, hourly\_wage) }\SpecialCharTok{\%\textgreater{}\%}

  \FunctionTok{drop\_na}\NormalTok{(number\_of\_years\_worked\_in\_the\_program, hourly\_wage) }\SpecialCharTok{\%\textgreater{}\%}

  \FunctionTok{group\_by}\NormalTok{(number\_of\_years\_worked\_in\_the\_program) }\SpecialCharTok{\%\textgreater{}\%}

  \FunctionTok{summarize}\NormalTok{(}

    \AttributeTok{n =} \FunctionTok{n}\NormalTok{(),}

    \AttributeTok{mean\_wage =} \FunctionTok{mean}\NormalTok{(hourly\_wage),}

    \AttributeTok{median\_wage =} \FunctionTok{median}\NormalTok{(hourly\_wage),}

    \AttributeTok{sd\_wage =} \FunctionTok{sd}\NormalTok{(hourly\_wage),}

    \AttributeTok{min\_wage =} \FunctionTok{min}\NormalTok{(hourly\_wage),}

    \AttributeTok{max\_wage =} \FunctionTok{max}\NormalTok{(hourly\_wage)}

\NormalTok{  )}

\NormalTok{wage\_by\_years}
\end{Highlighting}
\end{Shaded}

\begin{verbatim}
## # A tibble: 16 x 7
##    number_of_years_worke~1     n mean_wage median_wage sd_wage min_wage max_wage
##    <fct>                   <int>     <dbl>       <dbl>   <dbl>    <dbl>    <dbl>
##  1 (01) Less than half a ~   444      11.7        10.4    6.29     3.85    108. 
##  2 (02) 7-12 months          379      14.5        12     10.4      1.29    108. 
##  3 (03) 1 year               619      14.0        12      9.27     1.5     108. 
##  4 (04) 2 years              470      14.8        12.4    9.75     1.56    108. 
##  5 (05) 3 years              397      15.3        12.5   12.0      2.85    108. 
##  6 (06) 4 years              284      15.0        13     10.1      2.22    108. 
##  7 (07) 5 years              248      15.0        13.5    6.89     4.2      58.3
##  8 (08) 6 years              173      16.7        14     12.2      2.31    108. 
##  9 (09) 7 years              139      15.7        13.7    8.29     1.88     55.0
## 10 (10) 8 years              109      16.1        14.3    8.42     2.08     51.4
## 11 (11) 9 years               79      16.9        15      7.95     7.69     51.4
## 12 (12) 10 years             148      15.5        13.7    8.03     5.77     63.3
## 13 (13) 11 to 15 years       350      19.1        14.7   16.2      2.86    108. 
## 14 (14) 16 to 20 years       219      18.0        15.4   11.8      3.75    108. 
## 15 (15) 21 to 25 years       108      19.2        16.4   12.2      4.32    108. 
## 16 (16) 26 or more years      89      18.8        16     13.3      1.86    108. 
## # i abbreviated name: 1: number_of_years_worked_in_the_program
\end{verbatim}

The correlation coefficient summarizes the strength and direction of the
relationship between years worked in the program and hourly wage.

\begin{Shaded}
\begin{Highlighting}[]
\NormalTok{wage\_by\_years\_plot }\OtherTok{\textless{}{-}}
\NormalTok{  nsece\_2019\_wf\_recode }\SpecialCharTok{\%\textgreater{}\%}
  \FunctionTok{drop\_na}\NormalTok{(number\_of\_years\_worked\_in\_the\_program, hourly\_wage) }\SpecialCharTok{\%\textgreater{}\%}
  \FunctionTok{ggplot}\NormalTok{(}\FunctionTok{aes}\NormalTok{(}\AttributeTok{x =}\NormalTok{ number\_of\_years\_worked\_in\_the\_program,}
             \AttributeTok{y =}\NormalTok{ hourly\_wage)) }\SpecialCharTok{+}
  \FunctionTok{geom\_point}\NormalTok{() }\SpecialCharTok{+}
  \FunctionTok{geom\_smooth}\NormalTok{(}\AttributeTok{method =} \StringTok{"lm"}\NormalTok{)}

\NormalTok{wage\_by\_years\_plot}
\end{Highlighting}
\end{Shaded}

\begin{verbatim}
## `geom_smooth()` using formula = 'y ~ x'
\end{verbatim}

\pandocbounded{\includegraphics[keepaspectratio]{ProblemSet_1_files/ProblemSet_1_files/figure-latex/unnamed-chunk-11-1.pdf}}

The scatterplot shows whether hourly wage tends to increase, decrease,
or remain relatively constant as the number of years worked in the
program increases.

\begin{center}\rule{0.5\linewidth}{0.5pt}\end{center}

\subsection{c.~Hourly wage and
race/ethnicity}\label{c.-hourly-wage-and-raceethnicity}

Take \textbf{hourly\_wage} and \textbf{race/ethnicity} for example. We
can start by asking first; among the respondents, does hourly wage
differ across racial and ethnic groups? We can graph this relationship
using a boxplot.

\begin{Shaded}
\begin{Highlighting}[]
\CommentTok{\# c. Hourly wage and race}

\CommentTok{\# Research question:}

\CommentTok{\# Does hourly wage differ by race among ECE workers?}

\NormalTok{wage\_by\_race }\OtherTok{\textless{}{-}}

\NormalTok{  nsece\_2019\_wf\_recode }\SpecialCharTok{\%\textgreater{}\%}

  \FunctionTok{select}\NormalTok{(race, hourly\_wage) }\SpecialCharTok{\%\textgreater{}\%}

  \FunctionTok{drop\_na}\NormalTok{(race, hourly\_wage) }\SpecialCharTok{\%\textgreater{}\%}

  \FunctionTok{group\_by}\NormalTok{(race) }\SpecialCharTok{\%\textgreater{}\%}

  \FunctionTok{summarize}\NormalTok{(}

    \AttributeTok{n =} \FunctionTok{n}\NormalTok{(),}

    \AttributeTok{mean\_wage =} \FunctionTok{mean}\NormalTok{(hourly\_wage),}

    \AttributeTok{median\_wage =} \FunctionTok{median}\NormalTok{(hourly\_wage),}

    \AttributeTok{sd\_wage =} \FunctionTok{sd}\NormalTok{(hourly\_wage),}

    \AttributeTok{min\_wage =} \FunctionTok{min}\NormalTok{(hourly\_wage),}

    \AttributeTok{max\_wage =} \FunctionTok{max}\NormalTok{(hourly\_wage)}

\NormalTok{  )}

\NormalTok{wage\_by\_race}
\end{Highlighting}
\end{Shaded}

\begin{verbatim}
## # A tibble: 4 x 7
##   race                         n mean_wage median_wage sd_wage min_wage max_wage
##   <fct>                    <int>     <dbl>       <dbl>   <dbl>    <dbl>    <dbl>
## 1 (1) White Only            2607      15.5        13.0   10.3      1.5      108.
## 2 (2) Black or African Am~   882      13.7        11.5   10.0      1.29     108.
## 3 (3) Asian Only             186      17.5        16      9.68     3.31     107.
## 4 (8) Other                  204      15.9        13     11.6      5.71     108.
\end{verbatim}

The summary statistics show whether average hourly wage differs across
racial and ethnic groups.

\begin{Shaded}
\begin{Highlighting}[]
\NormalTok{wage\_by\_race\_plot }\OtherTok{\textless{}{-}}
\NormalTok{  nsece\_2019\_wf\_recode }\SpecialCharTok{\%\textgreater{}\%}
  \FunctionTok{drop\_na}\NormalTok{(race, hourly\_wage) }\SpecialCharTok{\%\textgreater{}\%}
  \FunctionTok{ggplot}\NormalTok{(}\FunctionTok{aes}\NormalTok{(}\AttributeTok{x =}\NormalTok{ race, }\AttributeTok{y =}\NormalTok{ hourly\_wage)) }\SpecialCharTok{+}
  \FunctionTok{geom\_boxplot}\NormalTok{()}

\NormalTok{wage\_by\_race\_plot}
\end{Highlighting}
\end{Shaded}

\pandocbounded{\includegraphics[keepaspectratio]{ProblemSet_1_files/ProblemSet_1_files/figure-latex/unnamed-chunk-13-1.pdf}}

The boxplot illustrates the distribution of hourly wages within each
racial and ethnic group, allowing differences in wages to be compared
visually.

\begin{center}\rule{0.5\linewidth}{0.5pt}\end{center}

\subsection{d.~Hourly wage and educational
attainment}\label{d.-hourly-wage-and-educational-attainment}

Take \textbf{hourly\_wage} and
\textbf{highest\_level\_of\_education\_completed} for example. We can
start by asking first; among the respondents, does hourly wage differ by
educational attainment? A boxplot can help visualize this relationship.

\begin{Shaded}
\begin{Highlighting}[]
\CommentTok{\# d. Hourly wage and educational attainment}

\CommentTok{\# Research question:}
\CommentTok{\# Does hourly wage differ by educational attainment among ECE workers?}

\NormalTok{wage\_by\_education }\OtherTok{\textless{}{-}}
\NormalTok{  nsece\_2019\_wf\_recode }\SpecialCharTok{\%\textgreater{}\%}
  \FunctionTok{select}\NormalTok{(highest\_level\_of\_education\_completed, hourly\_wage) }\SpecialCharTok{\%\textgreater{}\%}
  \FunctionTok{drop\_na}\NormalTok{(highest\_level\_of\_education\_completed, hourly\_wage) }\SpecialCharTok{\%\textgreater{}\%}
  \FunctionTok{group\_by}\NormalTok{(highest\_level\_of\_education\_completed) }\SpecialCharTok{\%\textgreater{}\%}
  \FunctionTok{summarize}\NormalTok{(}
    \AttributeTok{n =} \FunctionTok{n}\NormalTok{(),}
    \AttributeTok{mean\_wage =} \FunctionTok{mean}\NormalTok{(hourly\_wage),}
    \AttributeTok{median\_wage =} \FunctionTok{median}\NormalTok{(hourly\_wage),}
    \AttributeTok{sd\_wage =} \FunctionTok{sd}\NormalTok{(hourly\_wage),}
    \AttributeTok{min\_wage =} \FunctionTok{min}\NormalTok{(hourly\_wage),}
    \AttributeTok{max\_wage =} \FunctionTok{max}\NormalTok{(hourly\_wage)}
\NormalTok{  )}

\NormalTok{wage\_by\_education}
\end{Highlighting}
\end{Shaded}

\begin{verbatim}
## # A tibble: 7 x 7
##   highest_level_of_educa~1     n mean_wage median_wage sd_wage min_wage max_wage
##   <fct>                    <int>     <dbl>       <dbl>   <dbl>    <dbl>    <dbl>
## 1 (1) Less than High Scho~    75      10.3        9.75    3.91     4.44     28.3
## 2 (3) GED or high school ~   120      11.4       10.3     4.89     4.44     51.4
## 3 (4) High school graduate   710      11.5       10.5     5.51     1.29     51.4
## 4 (5) Some college credit~  1179      12.5       12       5.03     1.56     51.4
## 5 (6) Associate degree (A~   753      15.5       13.8    10.3      1.5     107. 
## 6 (7) Bachelor's degree (~  1052      18.6       15.5    13.1      2.5     108. 
## 7 (8) Graduate or profess~   374      24.1       20      16.7      1.86    108. 
## # i abbreviated name: 1: highest_level_of_education_completed
\end{verbatim}

The summary statistics show whether average hourly wage differs across
levels of educational attainment.

\begin{Shaded}
\begin{Highlighting}[]
\NormalTok{wage\_by\_education\_plot }\OtherTok{\textless{}{-}}
\NormalTok{  nsece\_2019\_wf\_recode }\SpecialCharTok{\%\textgreater{}\%}
  \FunctionTok{drop\_na}\NormalTok{(highest\_level\_of\_education\_completed, hourly\_wage) }\SpecialCharTok{\%\textgreater{}\%}
  \FunctionTok{ggplot}\NormalTok{(}\FunctionTok{aes}\NormalTok{(}\AttributeTok{x =}\NormalTok{ highest\_level\_of\_education\_completed, }\AttributeTok{y =}\NormalTok{ hourly\_wage)) }\SpecialCharTok{+}
  \FunctionTok{geom\_boxplot}\NormalTok{()}

\NormalTok{wage\_by\_education\_plot}
\end{Highlighting}
\end{Shaded}

\pandocbounded{\includegraphics[keepaspectratio]{ProblemSet_1_files/ProblemSet_1_files/figure-latex/unnamed-chunk-15-1.pdf}}

The boxplot shows the distribution of hourly wages across levels of
educational attainment, making it easier to compare wages among
respondents with different levels of education.

\begin{center}\rule{0.5\linewidth}{0.5pt}\end{center}

\subsection{e. Hourly wage and gender}\label{e.-hourly-wage-and-gender}

Take \textbf{hourly\_wage} and \textbf{gender} for example. We can start
by asking first; among the respondents, does hourly wage differ between
men and women? We can graph this relationship using a boxplot.

\begin{Shaded}
\begin{Highlighting}[]
\CommentTok{\# e. Hourly wage and gender}

\CommentTok{\# Research question:}
\CommentTok{\# Does hourly wage differ by gender among ECE workers?}

\NormalTok{wage\_by\_gender }\OtherTok{\textless{}{-}}
\NormalTok{  nsece\_2019\_wf\_recode }\SpecialCharTok{\%\textgreater{}\%}
  \FunctionTok{select}\NormalTok{(gender, hourly\_wage) }\SpecialCharTok{\%\textgreater{}\%}
  \FunctionTok{drop\_na}\NormalTok{(gender, hourly\_wage) }\SpecialCharTok{\%\textgreater{}\%}
  \FunctionTok{group\_by}\NormalTok{(gender) }\SpecialCharTok{\%\textgreater{}\%}
  \FunctionTok{summarize}\NormalTok{(}
    \AttributeTok{n =} \FunctionTok{n}\NormalTok{(),}
    \AttributeTok{mean\_wage =} \FunctionTok{mean}\NormalTok{(hourly\_wage),}
    \AttributeTok{median\_wage =} \FunctionTok{median}\NormalTok{(hourly\_wage),}
    \AttributeTok{sd\_wage =} \FunctionTok{sd}\NormalTok{(hourly\_wage),}
    \AttributeTok{min\_wage =} \FunctionTok{min}\NormalTok{(hourly\_wage),}
    \AttributeTok{max\_wage =} \FunctionTok{max}\NormalTok{(hourly\_wage)}
\NormalTok{  )}

\NormalTok{wage\_by\_gender}
\end{Highlighting}
\end{Shaded}

\begin{verbatim}
## # A tibble: 2 x 7
##   gender         n mean_wage median_wage sd_wage min_wage max_wage
##   <fct>      <int>     <dbl>       <dbl>   <dbl>    <dbl>    <dbl>
## 1 (1) Male      96      18.8        13.5    18.0     3.09     107.
## 2 (2) Female  4133      15.3        12.8    10.4     1.29     108.
\end{verbatim}

The summary statistics show whether average hourly wage differs between
men and women.

\begin{Shaded}
\begin{Highlighting}[]
\NormalTok{wage\_by\_gender\_plot }\OtherTok{\textless{}{-}}
\NormalTok{  nsece\_2019\_wf\_recode }\SpecialCharTok{\%\textgreater{}\%}
  \FunctionTok{drop\_na}\NormalTok{(gender, hourly\_wage) }\SpecialCharTok{\%\textgreater{}\%}
  \FunctionTok{ggplot}\NormalTok{(}\FunctionTok{aes}\NormalTok{(}\AttributeTok{x =}\NormalTok{ gender, }\AttributeTok{y =}\NormalTok{ hourly\_wage)) }\SpecialCharTok{+}
  \FunctionTok{geom\_boxplot}\NormalTok{()}

\NormalTok{wage\_by\_gender\_plot}
\end{Highlighting}
\end{Shaded}

\pandocbounded{\includegraphics[keepaspectratio]{ProblemSet_1_files/ProblemSet_1_files/figure-latex/unnamed-chunk-17-1.pdf}}

The boxplot shows the distribution of hourly wages for men and women,
making it easier to compare differences in wages between the two groups.

Citations Pacheco-Applegate, A. (2026). Data for Policy Analysis and
Management: Class 4---Bivariate descriptive statistics and basic
statistical tests {[}PowerPoint slides{]}. University of Chicago.

\begin{enumerate}
\def\labelenumi{\arabic{enumi}.}
\setcounter{enumi}{2}
\tightlist
\item
  Policy implications. Based on your analyses, describe the segment(s)
  of the ECE workforce that appear to be most affected by low wages.
  What characteristics would you consider when designing a policy to
  improve compensation in the ECE sector? Support your answer using
  evidence from your analyses. (2 points)
\end{enumerate}

Based on my analyses, aides and assistant teachers appear to be the
segment of the ECE workforce most affected by low wages, as they had the
lowest average hourly wage compared with teachers, instructors, and lead
teachers. The analyses also suggest that hourly wages differ across
race, educational attainment, gender, and years worked in the program,
indicating that compensation varies across multiple workforce
characteristics. When designing a policy to improve compensation in the
ECE sector, I would prioritize raising wages for lower-paid staff roles
while also considering differences in education, experience, and
demographic characteristics. Targeting these factors could help reduce
wage disparities and improve equity across the ECE workforce.

\end{document}
